% !TEX program = xelatex
% Самостоятельная работа по дисциплине «Основы финансовой математики»
% Тема: «Порядок оценки доходности финансового инструмента»
% Документ самодостаточный: компилируется XeLaTeX без внешних файлов.
\documentclass[14pt, a4paper]{extarticle}

%% Шрифты и язык
\usepackage{fontspec}
\setmainfont{Times New Roman}
\setmonofont{Courier New}

\usepackage{polyglossia}
\setmainlanguage{russian}
\setotherlanguage{english}

\newfontfamily\cyrillicfont{Times New Roman}
\newfontfamily\cyrillicfonttt{Courier New}

%% Поля страницы
\usepackage[
  a4paper,
  top=2cm,
  bottom=2cm,
  left=3cm,
  right=1.5cm
]{geometry}

%% Межстрочный интервал 1.5
\usepackage{setspace}
\onehalfspacing

%% Отступ первой строки абзаца
\usepackage{indentfirst}
\setlength{\parindent}{1.25cm}

%% Заголовки разделов
\usepackage{titlesec}
\titleformat{\section}
  {\normalfont\large\bfseries\centering}
  {\thesection}{0.5em}{}
\titlespacing*{\section}{0pt}{14pt}{8pt}

\titleformat{\subsection}
  {\normalfont\normalsize\bfseries}
  {\thesubsection}{0.5em}{}
\titlespacing*{\subsection}{0pt}{10pt}{6pt}

%% Математика, графика, таблицы
\usepackage{amsmath, amssymb}
\usepackage{graphicx}
\usepackage{booktabs}
\usepackage{array}

%% Гиперссылки
\usepackage[hidelinks]{hyperref}
\hypersetup{
  colorlinks=true,
  linkcolor=black,
  citecolor=black,
  urlcolor=cyan
}

%% Нумерация страниц — снизу по центру
\usepackage{fancyhdr}
\pagestyle{fancy}
\fancyhf{}
\fancyfoot[C]{\thepage}
\renewcommand{\headrulewidth}{0pt}

% -------------------------------------------------------
\begin{document}

%% Содержание
\renewcommand{\contentsname}{Содержание}
\begingroup
\normalsize
\tableofcontents
\endgroup
\newpage

% =======================================================
\section*{Введение}
\addcontentsline{toc}{section}{Введение}
\setcounter{section}{0}

Современный финансовый рынок предлагает инвестору широкий набор инструментов
для размещения свободных денежных средств: акции, облигации, банковские
депозиты, векселя, паи инвестиционных фондов, производные финансовые
инструменты. Каждый из них характеризуется собственным сочетанием доходности,
риска и ликвидности, а рациональный выбор возможен только при наличии единой,
сопоставимой меры эффективности вложений. Такой мерой выступает
\emph{доходность} финансового инструмента.

\textbf{Понятие финансовых инструментов.} Финансовый инструмент — это
договор, в результате которого у одной стороны возникает финансовый актив, а у
другой — финансовое обязательство или долевой инструмент. К финансовым
инструментам относят как первичные инструменты (денежные средства,
дебиторская и кредиторская задолженность, долевые и долговые ценные бумаги),
так и производные (фьючерсы, форварды, опционы, свопы). С точки зрения
инвестора финансовый инструмент — это способ вложения капитала с целью
получения будущего дохода в виде процентов, дивидендов, купонов или прироста
рыночной стоимости.

\textbf{Значение оценки доходности для инвесторов.} Оценка доходности
позволяет решить три ключевые задачи. Во-первых, измерить фактическую
эффективность уже совершённых вложений (ретроспективный анализ). Во-вторых,
спрогнозировать ожидаемый результат по планируемым вложениям (перспективный
анализ). В-третьих, обеспечить \emph{сопоставимость} разнородных инструментов,
приведя их к общей размерности — как правило, к годовой процентной ставке.
Без корректно рассчитанной доходности инвестор не может ни сравнить
альтернативы, ни оценить, компенсирует ли ожидаемый доход принимаемый риск.

Цель настоящей работы — систематизировать понятие доходности финансового
инструмента, рассмотреть её основные виды и показатели, а также описать
последовательный порядок (алгоритм) оценки доходности с иллюстрацией на
практическом примере.

% =======================================================
\section{Понятие доходности финансового инструмента}

\subsection{Что такое доходность}

\textbf{Доходность} — это относительный показатель, характеризующий величину
дохода, полученного (или ожидаемого к получению) на единицу вложенного
капитала за определённый период времени. В отличие от абсолютного дохода,
выражаемого в денежных единицах, доходность выражается в долях или процентах и
потому не зависит от масштаба вложений. Именно относительная природа делает
доходность универсальным измерителем, пригодным для сравнения инструментов
разной стоимости и разного срока.

Формально доходность за период определяется как отношение суммарного дохода к
величине первоначально вложенного капитала:
\begin{equation}
  r = \frac{D}{P_0},
\end{equation}
где $D$ — совокупный доход за период, $P_0$ — сумма первоначальных вложений
(цена приобретения инструмента).

\subsection{Виды доходности}

В финансовой практике различают несколько взаимодополняющих видов доходности.

\textbf{Абсолютная доходность} — это доход, выраженный в денежных единицах
(рублях, долларах). Строго говоря, это абсолютный \emph{доход}, а не
доходность в относительном смысле; тем не менее термин широко используется для
обозначения итогового денежного результата операции.

\textbf{Относительная доходность} — доход, отнесённый к вложенному капиталу и
выраженный в процентах. Это базовая форма представления доходности, дающая
возможность сравнивать инструменты между собой.

\textbf{Текущая доходность} (current yield) характеризует регулярный
периодический доход (купон, дивиденд) по отношению к текущей рыночной цене
инструмента, без учёта прироста или потери капитала:
\begin{equation}
  r_{\text{тек}} = \frac{C}{P},
\end{equation}
где $C$ — годовой купонный (дивидендный) доход, $P$ — текущая цена. Показатель
удобен для быстрой оценки, но игнорирует изменение стоимости тела инструмента.

\textbf{Ожидаемая доходность} (expected return) — прогнозная величина,
рассчитываемая как средневзвешенное возможных исходов с учётом вероятностей их
наступления:
\begin{equation}
  E(r) = \sum_{i=1}^{n} p_i \, r_i, \qquad \sum_{i=1}^{n} p_i = 1,
\end{equation}
где $r_i$ — доходность в $i$-м сценарии, $p_i$ — вероятность этого сценария.
Ожидаемая доходность лежит в основе портфельной теории и оценки риска.

\textbf{Эффективная доходность} (effective yield) учитывает реинвестирование
промежуточных доходов и капитализацию процентов, приводя все денежные потоки к
единой годовой базе. Именно эффективная доходность корректно отражает реальную
отдачу от инструмента и используется для сопоставления вложений с разной
периодичностью начисления.

% =======================================================
\section{Основные показатели оценки доходности}

\subsection{Доход (прибыль) от финансового инструмента}

Совокупный доход от владения финансовым инструментом складывается из двух
компонентов: \emph{текущего дохода} (проценты, купоны, дивиденды) и
\emph{дохода от прироста капитала} (разница между ценой продажи и ценой
покупки). В общем виде:
\begin{equation}
  D = (P_1 - P_0) + CF,
\end{equation}
где $P_0$ — цена покупки, $P_1$ — цена продажи (или текущая стоимость), $CF$ —
сумма полученных за период процентных, купонных или дивидендных выплат. Прирост
капитала может быть как положительным, так и отрицательным.

\subsection{Норма (ставка) доходности}

Норма доходности — это доход, выраженный в процентах годовых. Она позволяет
свести разные по срокам операции к единой временной базе и является главным
критерием при сравнении инструментов. Различают \emph{номинальную} ставку
(объявленную, без учёта капитализации и инфляции) и \emph{реальную} ставку
(очищенную от инфляции).

\subsection{Простая и сложная процентные ставки}

При начислении по \textbf{простой} процентной ставке проценты начисляются
только на первоначальную сумму и не капитализируются. Наращённая сумма:
\begin{equation}
  FV = PV\,(1 + i \cdot n),
\end{equation}
где $PV$ — первоначальная сумма, $i$ — ставка за период, $n$ — число периодов.

При начислении по \textbf{сложной} процентной ставке проценты присоединяются к
основной сумме и в следующем периоде сами приносят доход (капитализация):
\begin{equation}
  FV = PV\,(1 + i)^{n}.
\end{equation}
Именно эффект сложного процента обеспечивает опережающий рост капитала на
длинных горизонтах инвестирования.

\subsection{Годовая доходность (Annual Yield)}

Если доходность получена за период, отличный от года, её приводят к годовой
базе. При простом (пропорциональном) пересчёте:
\begin{equation}
  r_{\text{год}} = r \cdot \frac{T}{t},
\end{equation}
где $r$ — доходность за фактический срок владения, $t$ — число дней (месяцев)
владения, $T$ — число дней (месяцев) в году. При учёте капитализации
используется \textbf{эффективная годовая ставка} (Effective Annual Rate):
\begin{equation}
  i_{\text{eff}} = \left(1 + \frac{j}{m}\right)^{m} - 1,
\end{equation}
где $j$ — номинальная годовая ставка, $m$ — число начислений процентов в году.
Эффективная ставка всегда не меньше номинальной и совпадает с ней лишь при
$m = 1$.

% =======================================================
\section{Порядок оценки доходности}

Оценка доходности финансового инструмента представляет собой
последовательность логически связанных этапов.

\subsection{Сбор исходных данных}

На первом этапе фиксируются все параметры операции:
\begin{itemize}
  \item \textbf{цена покупки} ($P_0$) — сумма, фактически уплаченная за
        инструмент, включая комиссии и издержки приобретения;
  \item \textbf{цена продажи или текущая стоимость} ($P_1$) — сумма, за
        которую инструмент реализован либо оценивается на текущий момент;
  \item \textbf{полученные выплаты} ($CF$) — проценты по вкладу, купоны по
        облигации, дивиденды по акции за период владения;
  \item \textbf{срок владения} ($t$) — период между приобретением и продажей
        (оценкой), выраженный в днях, месяцах или годах.
\end{itemize}

\subsection{Выбор метода расчёта}

Метод расчёта определяется характером инструмента и целью анализа. Для
краткосрочных операций без реинвестирования применяют простую доходность; для
инструментов с капитализацией и промежуточными выплатами — эффективную
доходность или показатель внутренней нормы доходности (IRR). Для прогнозных
оценок используют формулу ожидаемой доходности.

\subsection{Расчёт доходности}

На этом этапе выбранные формулы наполняются исходными данными и выполняется
собственно вычисление: определяется абсолютный доход, затем простая доходность
за срок владения и, при необходимости, её приведение к годовой (эффективной)
базе.

\subsection{Анализ полученного результата}

Полученное значение интерпретируется: оценивается его знак и величина,
сравнивается с ожиданиями инвестора, а также с уровнем инфляции для оценки
реальной (а не только номинальной) отдачи.

\subsection{Сравнение с альтернативными инструментами}

Заключительный этап — сопоставление рассчитанной годовой доходности с
доходностью альтернативных вложений сопоставимого уровня риска (например, со
ставкой по банковскому депозиту или доходностью государственных облигаций). Это
позволяет принять обоснованное инвестиционное решение.

% =======================================================
\section{Основные формулы}

Обобщим ключевые расчётные соотношения, используемые при оценке доходности.

\textbf{Абсолютный доход:}
\begin{equation}
  D = (P_1 - P_0) + CF.
\end{equation}

\textbf{Простая доходность за период владения:}
\begin{equation}
  r = \frac{(P_1 - P_0) + CF}{P_0}.
\end{equation}

\textbf{Доходность с учётом времени (годовая):}
\begin{equation}
  r_{\text{год}} = \frac{(P_1 - P_0) + CF}{P_0} \cdot \frac{T}{t}.
\end{equation}

\textbf{Эффективная годовая доходность (с учётом капитализации):}
\begin{equation}
  i_{\text{eff}} = \left(1 + r\right)^{\frac{T}{t}} - 1.
\end{equation}

% =======================================================
\section{Практический пример}

\subsection{Исходные данные}

Инвестор приобрёл корпоративную облигацию по цене $P_0 = 1000$ руб. Через
$t = 180$ дней облигация была продана по цене $P_1 = 1030$ руб. За период
владения инвестор получил купонный доход $CF = 40$ руб. Число дней в году
принимается равным $T = 365$.

\subsection{Пошаговый расчёт}

\textbf{Шаг 1. Абсолютный доход.}
\begin{equation*}
  D = (P_1 - P_0) + CF = (1030 - 1000) + 40 = 70 \text{ руб.}
\end{equation*}

\textbf{Шаг 2. Простая доходность за срок владения.}
\begin{equation*}
  r = \frac{D}{P_0} = \frac{70}{1000} = 0{,}07 = 7{,}0\,\%.
\end{equation*}

\textbf{Шаг 3. Приведение к годовой доходности.}
\begin{equation*}
  r_{\text{год}} = r \cdot \frac{T}{t}
  = 0{,}07 \cdot \frac{365}{180} \approx 0{,}1420 = 14{,}2\,\%.
\end{equation*}

\textbf{Шаг 4. Эффективная годовая доходность (с учётом капитализации).}
\begin{equation*}
  i_{\text{eff}} = (1 + r)^{\frac{T}{t}} - 1
  = (1{,}07)^{\frac{365}{180}} - 1 \approx 0{,}1470 = 14{,}7\,\%.
\end{equation*}

\subsection{Интерпретация результата}

За 180 дней инструмент принёс инвестору 7,0\,\% дохода, что в пересчёте на
год соответствует примерно 14,2\,\% простой годовой доходности и 14,7\,\%
эффективной доходности с учётом капитализации. Если ставка по сопоставимому по
риску банковскому депозиту составляет, например, 12\,\% годовых, то
рассмотренная облигация оказывается более привлекательным вложением. При этом
следует помнить, что при уровне инфляции 8\,\% реальная доходность инструмента
составит около $14{,}2 - 8 = 6{,}2\,\%$.

Сведём результаты расчёта в таблицу~\ref{tab:calc}.

\begin{table}[h]
  \centering
  \caption{Результаты оценки доходности облигации}
  \label{tab:calc}
  \begin{tabular}{lc}
    \toprule
    \textbf{Показатель} & \textbf{Значение} \\
    \midrule
    Абсолютный доход, руб. & 70 \\
    Простая доходность за 180 дней, \% & 7{,}0 \\
    Простая годовая доходность, \% & 14{,}2 \\
    Эффективная годовая доходность, \% & 14{,}7 \\
    Реальная доходность (при инфляции 8\,\%), \% & 6{,}2 \\
    \bottomrule
  \end{tabular}
\end{table}

% =======================================================
\section{Факторы, влияющие на доходность}

Фактическая и ожидаемая доходность финансового инструмента формируется под
воздействием ряда факторов.

\begin{itemize}
  \item \textbf{Рыночная конъюнктура.} Соотношение спроса и предложения,
        фаза экономического цикла и настроения инвесторов напрямую влияют на
        рыночные цены и, следовательно, на доходность от прироста капитала.
  \item \textbf{Инфляция.} Обесценивая денежные потоки, инфляция снижает
        реальную доходность; при высокой инфляции номинально прибыльное
        вложение может оказаться убыточным в реальном выражении.
  \item \textbf{Процентные ставки.} Рост ключевой ставки повышает доходность
        новых долговых инструментов, но снижает рыночную цену ранее выпущенных
        облигаций с фиксированным купоном (обратная зависимость цены и ставки).
  \item \textbf{Кредитный риск.} Чем выше вероятность дефолта эмитента, тем
        большую премию за риск требует инвестор, что выражается в повышенной
        требуемой доходности.
  \item \textbf{Ликвидность.} Инструменты, которые трудно быстро продать без
        потери в цене, несут премию за ликвидность и, как правило, должны
        обеспечивать более высокую доходность.
  \item \textbf{Срок инвестирования.} Более длинные вложения обычно требуют
        более высокой доходности как компенсации за возросшую неопределённость
        и упущенные альтернативы (временная структура процентных ставок).
\end{itemize}

% =======================================================
\section{Преимущества и ограничения методов оценки доходности}

\textbf{Простота расчётов.} Базовые показатели доходности (простая и годовая
доходность) вычисляются по несложным формулам и требуют минимума исходных
данных, что делает их доступными широкому кругу инвесторов.

\textbf{Возможность сравнения различных инструментов.} Приведение доходности к
единой годовой базе позволяет сопоставлять инструменты, различающиеся по цене,
сроку и характеру денежных потоков, и на этой основе ранжировать альтернативы.

\textbf{Ограничения без учёта риска и инфляции.} Главный недостаток простых
показателей — они отражают лишь ожидаемую или фактическую отдачу, но не
учитывают \emph{риск} её недополучения. Два инструмента с одинаковой
доходностью могут существенно различаться по волатильности. Кроме того,
номинальная доходность не отражает потерю покупательной способности денег из-за
инфляции. Для более полной оценки применяют показатели, учитывающие риск
(коэффициент Шарпа, стандартное отклонение доходности), и переходят от
номинальной к реальной доходности.

% =======================================================
\section*{Заключение}
\addcontentsline{toc}{section}{Заключение}

В работе рассмотрен порядок оценки доходности финансового инструмента как
ключевого критерия принятия инвестиционных решений.

\textbf{Основные выводы.} Доходность — это относительный показатель
эффективности вложений, выражаемый в процентах и не зависящий от масштаба
инвестиций. Существует несколько взаимодополняющих видов доходности
(абсолютная, относительная, текущая, ожидаемая, эффективная), каждый из которых
решает свою аналитическую задачу. Корректная оценка опирается на разграничение
простой и сложной процентных ставок и обязательное приведение результата к
единой годовой (эффективной) базе. Порядок оценки включает пять этапов: сбор
исходных данных, выбор метода расчёта, вычисление доходности, анализ результата
и сравнение с альтернативами.

\textbf{Практическое значение.} Рассчитанная и корректно
интерпретированная доходность позволяет инвестору сопоставлять разнородные
инструменты, оценивать реальную (с поправкой на инфляцию) отдачу вложений и
принимать обоснованные решения о размещении капитала. При этом показатели
доходности следует использовать в связке с оценкой риска и ликвидности, что
обеспечивает всесторонний анализ инвестиционной привлекательности финансовых
инструментов.

\end{document}
